\textbf{Privacy and responsible analytics.} The deployed runtime does not identify people and does not infer
demographic attributes: it stores track boxes, counts, crossings and dwell, and the VLM is instructed not to
identify anyone. Video frames are processed in memory and not recorded unless an output is configured; events,
time series, scene reports and chat are deleted after a configurable retention period (72\,h by default), while the
audit log of configuration changes is kept for accountability. The API requires a token per user with a role, and
the server binds to the local network by default. The campus case study (Appendix~\ref{app:campus}) is an offline
research analysis: its gender estimate is inferred from appearance by a probe trained on labels that were
themselves inferred from appearance, it can misclassify people and does not represent gender identity, it is
reported only in aggregate, and it is not needed for any crowd-safety function, so it is not part of the deployed
system. Recording people in public or institutional spaces requires notice and a lawful basis under local rules.
\textbf{Data provenance.} The adaptation video aa1 is a publicly available pedestrian benchmark video that is
distributed without an explicit licence; we use it for research only and do not redistribute its frames. The
held-out video test1 was recorded by the authors on the university campus for this study.
% AUTHORS: add the provenance of the two campus videos (video_1.mp4, video_2.mp4 = scenes 1 and 2) here.
Only derived data (boxes, tracks and aggregate metrics) and code are released by default; frames are shared only
where the video's licence and consent basis allow it.
